\ifmostrarGabarito
\begin{minted}{python}
from typing import List, Dict, Any, Tuple

def processar_transacoes_bancarias(
    transacoes: List[Dict[str, Any]],
    saldo_inicial: float
) -> Tuple[float, List[int], List[str]]:
    saldo_atual = float(saldo_inicial)
    sucessos = []
    erros = []
    
    for t in transacoes:
        try:
            # Validação de presença das chaves
            if "id" not in t or "tipo" not in t or "valor" not in t:
                raise KeyError("Chave ausente")
                
            t_id = t["id"]
            tipo = t["tipo"]
            valor = float(t["valor"])
            
            # Validação do valor estritamente positivo
            if valor <= 0:
                raise ValueError(f"Transacao {t_id}: Valor invalido")
                
            # Processamento por tipo
            if tipo == "DEPOSITO":
                saldo_atual += valor
                sucessos.append(t_id)
            elif tipo == "SAQUE":
                if valor > saldo_atual:
                    raise ValueError(f"Transacao {t_id}: Saldo insuficiente")
                saldo_atual -= valor
                sucessos.append(t_id)
            else:
                raise ValueError(f"Transacao {t_id}: Tipo desconhecido")
                
        except KeyError:
            erros.append("Erro na transacao: Chave ausente")
        except ValueError as e:
            erros.append(str(e))
            
    return (saldo_atual, sucessos, erros)
\end{minted}

\textbf{Explicação:}
Envolvemos o processamento de cada transação em um bloco \texttt{try/except}, disparando exceções apropriadas (\texttt{KeyError} e \texttt{ValueError}) e capturando-as para alimentar a lista de erros sem interromper a execução do lote.
\else
\linhasResposta{17}
\clearpage
\linhasResposta[FOLHA DE CONTINUACAO / RASCUNHO -- QUESTAO Q14]{30}
\clearpage
\fi
